\subsection{有界测度}

设 X 是一个局部紧空间。由于 \(\mathcal{K}(X;\mathbf{C})\) 上由 \(\mathcal{C}^{b}(X;\mathbf{C})\) 的拓扑诱导的拓扑粗于 \(\mathcal{K}(X;\mathbf{C})\) 上的直接极限拓扑，X 上的测度未必对 X 中的一致收敛拓扑连续。

定义 3. --- 局部紧空间 X 上的测度称为有界的，如果它在 \(\mathcal{K}(\mathrm{X};\mathbf{C})\) 上对一致收敛拓扑连续。

这等价于说存在一个有限数 \(M \geqslant 0\) ，使得对每个函数 \(f \in \mathcal{K}(\mathbf{X}; \mathbf{C})\) ，

\[
| \mu (f) | \leqslant \mathrm{M} \| f \|\tag{19}
\]

（其中 \(\| f\|\) 由 No. 2 的公式 (3) 定义）。

因此，说 \(\mu\) 是一个有界测度，就是说 \(\mu\) 属于空间 \(\mathcal{K}(\mathrm{X};\mathbf{C})\) （赋以范数 \(\| f\|\) ）的对偶；我们将把这个对偶记作 \(\mathcal{M}^1 (\mathrm{X};\mathbf{C})\) （如不致混淆，简记作 \(\mathcal{M}^1 (\mathrm{X})\) ）。我们知道 \(\mathcal{M}^1 (\mathrm{X};\mathbf{C})\) 赋有一个范数，\(\| \mu \|\) 是数 \(M\geqslant 0\) ——使不等式 (19) 对每个函数 \(f\in \mathcal{K}(\mathrm{X};\mathbf{C})\) 成立的那些——中最小者，或者写作

\[
\| \mu \| = \sup _ {\| f \| \leqslant 1, f \in \mathcal{K} (\mathrm{X}; \mathbf {C})} | \mu (f) |.\tag{20}
\]

赋以这一范数后，已知 \(\mathcal{M}^{1}(X;\mathbf{C})\) 是一个 Banach 空间（TVS, III, §3, No. 8, 命题 12 的推论 2）。

由公式 (20) 给出的 \(\|\mu\|\) 的定义可以延拓到 X 上的每个测度 \(\mu\) ，并且滥用地仍称 \(\|\mu\|\) 为 \(\mu\) 的范数；为使 \(\mu\) 有界，必须且只须 \(\|\mu\|\) 有限。

若 X 是紧的，则 X 上的每个测度都是有界的。

例子. --- 1) 测度 \(\varepsilon_{a}\) （由一点 \(a \in X\) 处的单位质量定义的）是有界的，且 \(\| \varepsilon_{a} \| = 1\) 。

\begin{enumerate}
\def\labelenumi{\arabic{enumi})}
\setcounter{enumi}{1}
\tightlist
\item
  \(\mathbf{R}\) 上的 Lebesgue 测度不是有界的；事实上，对每个整数 \(n > 0\) ，存在一个函数 \(f \in \mathcal{K}(\mathbf{R}; \mathbf{C})\) ，取值于 {[}0,1{]} 且在区间 \([-n, n]\) 上等于 1（No. 2, 引理 1）；于是 \(\| f \| = 1\) 且
\end{enumerate}

\[
\int_ {- \infty} ^ {+ \infty} f (x) d x \geqslant \int_ {- n} ^ {n} f (x) d x = 2 n,
\]

这证明不存在满足关系 (19) 的任何有限数 \(M\)。

\begin{enumerate}
\def\labelenumi{\arabic{enumi})}
\setcounter{enumi}{2}
\tightlist
\item
  在实直线 \(\mathbf{R}\) 上，映射
\end{enumerate}

\[
f \mapsto \int_ {- \infty} ^ {+ \infty} \frac {f (x) d x}{1 + x ^ {2}}
\]

是一个有界测度，因为对每个函数 \(f \in \mathcal{K}(\mathbf{R};\mathbf{C})\) ，

\[
\left| \int_ {- \infty} ^ {+ \infty} \frac {f (x) d x}{1 + x ^ {2}} \right| \leqslant \| f \| \int_ {- \infty} ^ {+ \infty} \frac {d x}{1 + x ^ {2}} = \pi \cdot \| f \|.
\]

由于关系 \(\| f\| \leqslant 1\) 与 \(\| \overline{f}\| \leqslant 1\) 等价，由 (20) 推出

\[
\| \overline {{\mu}} \| = \| \mu \|\tag{21}
\]

对 X 上的每个测度 \(\mu\) 成立。

命题 10. --- 对 X 上的每个测度 \(\mu\) ，

\[
\| \mu \| = \sup _ {0 \leqslant f \leqslant 1, f \in \mathcal {K} (\mathrm{X}; \mathbf {R})} | \mu | (f).\tag{22}
\]

因为，考虑到定义测度绝对值的公式 (12)，(22) 的第二项可以写成

\[
\sup _ {0 \leqslant f \leqslant 1, f \in \mathcal {K} (\mathrm{X}; \mathbf {R})} \left(\sup _ {| g | \leqslant f, g \in \mathcal {K} (\mathrm{X}; \mathbf {C})} | \mu (g) |\right) = \sup _ {\| g \| \leqslant 1, g \in \mathcal {K} (\mathrm{X}; \mathbf {C})} | \mu (g) |.
\]

推论 1. --- 对 X 上的每个测度 \(\mu\) ，\(\mu\) 与 \(|\mu|\) 的范数相等；\(\mu\) 有界当且仅当 \(|\mu|\) 有界。

推论 2. --- 对每个测度 \(\mu\) （紧空间 \(X\) 上的），

\[
\| \mu \| = | \mu | (1) = \int d | \mu |.\tag{23}
\]

这一公式将在 Ch. IV, §4, No. 7 中得到推广。

在紧空间 X 上，对每个（复）测度 \(\mu\) （X 上的），复数 \(\mu(1)\) 称为 \(\mu\) 的总质量。当 \(\mu\) 为正时，其总质量等于它的范数。当 \(\mu\) 是紧空间 X 上总质量等于 1 的正测度时，也称它的值 \(\mu(f)\) （对连续函数 \(f \in \mathcal{C}(\mathrm{X};\mathbf{C})\) 所取的）为 f 关于测度 \(\mu\) 的平均（值）。

推论 3. --- 对每个实测度 \(\mu\) （局部紧空间 \(X\) 上的），

\[
\| \mu \| = \sup _ {\| f \| \leqslant 1, f \in \mathcal{K} (\mathrm{X}; \mathbf {R})} | \mu (f) |.\tag{24}
\]

只须利用公式 (22) 以及 \(|\mu|(f)\) ——当 \(\mu\) 是实测度且 \(f \in \mathcal{K}_{+}(\mathrm{X})\) 时——的表达式 (9)。

因此有界实测度所成的集是赋范空间 \(\mathcal{K}(\mathrm{X};\mathbf{R})\) 的对偶；记作 \(\mathcal{M}^{1}(\mathrm{X},\mathbf{R})\) ，如不致混淆也记作 \(\mathcal{M}^{1}(\mathrm{X})\) 。典范单射 \(\mathcal{M}^{1}(\mathrm{X},\mathbf{R})\to\mathcal{M}^{1}(\mathrm{X};\mathbf{C})\) 由 (24) 是一个等距。

命题 11. --- 若 \(\mu\) 与 \(\nu\) 是 X 上的两个正测度，则 \(\| \mu + \nu \| = \| \mu \| + \| \nu \|\) 。

因为，函数 \(f \in \mathcal{K}(\mathrm{X}; \mathbf{R})\) （满足 \(0 \leqslant f \leqslant 1\) 的那些）对关系 \(\leqslant\) 构成一个有向集 S。于是对正测度 \(\mu\) （X 上的），由 (22) 与单调极限定理可得 \(\| \mu \| = \lim_{f \in S} \mu(f)\) ；命题的结论随之立即得出。

推论 1. --- 若 \(\mu\) 与 \(\nu\) 是 X 上的两个正测度且 \(\mu \leqslant \nu\) ，则 \(\| \mu \| \leqslant \| \nu \|\) ；特别地，若 \(\nu\) 有界则 \(\mu\) 亦有界。事实上，\(\| \nu \| = \| \mu \| + \| \nu - \mu \|\) 。

推论 2. --- 对每个实测度 \(\mu\) （X 上的），

\[
\| \mu \| = \| \mu^ {+} \| + \| \mu^ {-} \|.
\]

因为（命题 10 的推论 1），\(\mu\) 的范数等于 \(|\mu| = \mu^{+} + \mu^{-}\) 的范数。

命题 12. --- 若 \(\mu\) 是 X 上的一个有界测度，且 \(g\) 是 X 到 C 中的一个有界连续映射，则测度 \(g \cdot \mu\) 有界且 \(\| g \cdot \mu \| \leqslant \| g \| \cdot \| \mu \|\) 。

对每个函数 \(f\in \mathcal{K}(\mathrm{X};\mathbf{C})\)

\[
| \mu (f g) | \leqslant \| \mu \| \cdot \| f g \| \leqslant \| \mu \| \cdot \| g \| \cdot \| f \|.
\]

